%% ============================================================
%% Section 4: Technical Countermeasures (V0.2)
%% Source: V0.1 section5.tex (was §5, now §4 after swap)
%% Target: ~3,000 words
%% Changes: §4.1 Detection compressed (-600), §4.2 Watermarking compressed (-450),
%%   §4.3.1 Alignment compressed (-100), §4.3.2 Defense preserved,
%%   §4.3.3 Agentic Security EXPANDED (+300, blue ocean),
%%   New Table 4 (Detection) and Table 5 (Watermarking)
%% ============================================================

\section{Technical Countermeasures}
\label{sec:countermeasures}

To counter the threats identified in Section~\ref{sec:threats}, a multi-layered technical defense ecosystem has emerged spanning detection, watermarking, alignment, and agentic security. No single technique is sufficient; resilient AI requires composing all layers into an integrated security posture.


%% ------------------------------------------------------------
\subsection{AIGC Detection}
\label{sec:detect}
%% ------------------------------------------------------------

\subsubsection{Text Detection}
\label{sec:detect-text}

Text detection methods fall into two paradigms. \emph{Zero-shot statistical methods} exploit probability geometry of LM-generated text without labeled training data: DetectGPT~\cite{mitchell2023detectgpt} uses perturbation-based curvature estimation, Fast-DetectGPT~\cite{bao2023fastdetectgpt} achieves orders-of-magnitude speedups, and Binoculars~\cite{hans2024binoculars} computes cross-model agreement signals for black-box settings. \emph{Trained classifiers} fine-tune encoders on labeled pairs: GPTZero~\cite{adam2026gptzero} explicitly models Human/AI/Mixed authorship with sentence-level localization, while RADAR~\cite{hu2023radar} addresses paraphrase evasion through adversarial training. Turnitin's deployment across 200+ million papers reports $<$1\% document-level false positives at $\geq$20\% AI writing~\cite{turnitin2024}. Table~\ref{tab:detection} provides a detailed comparison.

Key challenges remain severe: paraphrase attacks cause 18.4-point accuracy drops~\cite{dugan2024raid}; cross-model generalization is weak (95\%+ in-distribution, $<$60\% cross-generator); multilingual performance degrades sharply (AUC from 0.946 to 0.537, English to German)~\cite{macko2023multitude}; and human--AI hybrid text produces the highest false positive rates~\cite{turnitin2024}. Benchmarks including RAID~\cite{dugan2024raid} (6M+ generations, 11 attacks), MULTITuDE~\cite{macko2023multitude} (11 languages), and M4~\cite{wang2024m4} (temporal evaluation) provide standardized assessment.

\subsubsection{Image, Audio, and Video Detection}
\label{sec:detect-multimodal}

For images, the critical challenge is cross-generator transfer: ProGAN-trained classifiers achieve as low as 3.05\% accuracy on diffusion outputs~\cite{ojha2023univfd}. CLIP-based approaches (UnivFD~\cite{ojha2023univfd}) and diffusion reconstruction methods (DIRE~\cite{wang2023dire}, AEROBLADE~\cite{ricker2024aeroblade}) improve generalization but remain fragile under benign processing---ResNet-50 accuracy drops from 96.2\% to $\sim$51\% under JPEG compression (q=30)~\cite{zhu2023genimage}.

For audio, the ASVspoof challenge series~\cite{asvspoof2019,asvspoof2021} formalizes spoofed-speech detection across TTS, voice conversion, and replay conditions. For video, FaceForensics++~\cite{rossler2019faceforensics} remains canonical, but diffusion-based video generation (Sora-class) produces entirely generated frames where existing face-swap detectors fail to transfer~\cite{cai2024beyonddeepfake}, requiring explicitly spatiotemporal approaches~\cite{stall2026}.

\subsubsection{Cross-Modal Challenges}
\label{sec:detect-cross}

Three bottlenecks converge across modalities: \emph{generalization} failure on new architectures, \emph{robustness} collapse under benign processing, and \emph{false positive harm} from treating scores as authoritative. The arms race is structurally asymmetric: generators cheaply change parameters while detectors must remain robust across all perturbations. No unified cross-modal framework exists.

\begin{table}[htbp]
\centering
\caption{Representative AIGC detection methods across modalities.}
\label{tab:detection}
\small
\begin{tabular}{p{1.3cm}p{2.0cm}p{2.8cm}p{2.5cm}p{3.5cm}}
\toprule
\textbf{Modality} & \textbf{Paradigm} & \textbf{Methods} & \textbf{Mechanism} & \textbf{Key Limitation} \\
\midrule
Text & Zero-shot & DetectGPT, Binoculars, DNA-GPT & Probability curvature / cross-model agreement & Requires generator access or proxy; paraphrase fragile \\
\addlinespace
Text & Supervised & GPTZero, RADAR, Turnitin & Encoder fine-tuning; adversarial training & Overfits to training generators; staleness \\
\addlinespace
Image & Frequency & Spectral analysis & GAN upsampling artifacts & Mitigated by post-processing \\
\addlinespace
Image & Representation & UnivFD (CLIP), LASTED & Feature-space classification & GAN$\to$diffusion transfer fails (3.05\%) \\
\addlinespace
Image & Reconstruction & DIRE, AEROBLADE & Diffusion reconstruction error & Computationally expensive; format-bias risk \\
\addlinespace
Audio & Challenge-based & ASVspoof series & Spectral/prosodic/codec features & Codec artifacts confound interpretation \\
\addlinespace
Video & Temporal & FaceForensics++, rPPG & Frame artifacts; physiological signals & Sora-class generators: no face-swap artifacts \\
\bottomrule
\end{tabular}
\end{table}


%% ------------------------------------------------------------
\subsection{Watermarking and Content Provenance}
\label{sec:watermark}
%% ------------------------------------------------------------

\subsubsection{Text Watermarking}
\label{sec:wm-text}

The foundational KGW scheme~\cite{kirchenbauer2023kgw} biases token selection toward pseudorandom ``green list'' tokens, enabling statistical detection via z-test. Subsequent work improved robustness (Unigram watermark with fixed partitions~\cite{zhao2023unigram}), capacity (multi-bit encoding via Reed--Solomon codes~\cite{qu2024multibit}), and semantic resilience (SemStamp's sentence-level embedding hashing~\cite{hou2023semstamp}). A distinct paradigm achieves watermarking \emph{without altering the output distribution}: Christ et al.~\cite{christ2024undetectable} formalize cryptographic sampling where watermarked and original distributions are provably identical. Google DeepMind's SynthID~\cite{synthidtext2024} is the first production deployment.

The core tension is a three-way trade-off between quality, strength, and robustness. Paraphrasing via SIRA achieves $\sim$100\% watermark removal~\cite{cheng2025sira}; spoofing attacks forge signals in human text~\cite{sadasivan2023spoofing}. Cross-provider standardization remains absent. Table~\ref{tab:watermark} compares representative schemes.

\subsubsection{Image Watermarking}
\label{sec:wm-image}

Two paradigms dominate. \emph{Diffusion-native methods} embed watermarks during generation: Stable Signature~\cite{fernandez2023stablesig} in the decoder ($>$90\% accuracy with 90\% cropping), Tree-Ring~\cite{wen2023treering} in initial noise (survives JPEG and geometric transforms), and Gaussian Shading~\cite{yang2024gaussianshading} in latent variables (zero quality loss). \emph{Encoder--decoder methods} (HiDDeN~\cite{zhu2018hidden}, StegaStamp~\cite{tancik2020stegastamp}) apply watermarks post-generation with $\sim$95\% bit recovery after print-and-capture. Google's SynthID-Image, deployed on $>$10 billion images, achieves $>$99\% detection under common edits~\cite{gowal2024synthidimage}. Regeneration attacks can remove weak watermarks but degrade image quality for strong ones~\cite{zhao2023regeneration}.

\subsubsection{Content Provenance}
\label{sec:provenance}

The C2PA standard (v2.2, May 2025) embeds cryptographically signed ``Content Credentials'' with hardware adoption (Leica, Nikon, Google Pixel~10) and platform detection (Meta, Google, TikTok)~\cite{c2pa2025}. CISA endorsed Content Credentials in January 2025~\cite{cisa2025cc}. Limitations include metadata stripping on social media upload and the chicken-and-egg adoption problem. The emerging direction combines credentials with invisible watermarks to survive metadata loss.

\subsubsection{Challenges}
\label{sec:wm-challenges}

Christ et al.\ prove perfectly undetectable watermarks require cryptographic secrecy~\cite{christ2024undetectable}. Open-source model watermarks can be removed by weight merging (AUC $\to$ 0.5)~\cite{xu2025removal}. Cross-provider standardization is absent. No watermark is secure against a dedicated white-box adversary.

\begin{table}[htbp]
\centering
\caption{Representative watermarking schemes for AI-generated content.}
\label{tab:watermark}
\small
\begin{tabular}{p{1.2cm}p{2.2cm}p{2.0cm}p{1.5cm}p{1.5cm}p{3.5cm}}
\toprule
\textbf{Modality} & \textbf{Method} & \textbf{Embedding} & \textbf{Capacity} & \textbf{Quality Loss} & \textbf{Robustness} \\
\midrule
Text & KGW & Sampling bias & 1 bit & Moderate & Fragile to paraphrase \\
\addlinespace
Text & Christ et al. & Cryptographic sampling & 1 bit & Zero & Survives $\sim$40\% edits \\
\addlinespace
Text & SynthID & Token probability & Multi-bit & Low & Moderate (long text) \\
\addlinespace
Image & Stable Signature & Decoder fine-tuning & Multi-bit & Low & $>$90\% at 90\% crop \\
\addlinespace
Image & Tree-Ring & Initial noise & 1 bit & Zero & JPEG, rotation, crop \\
\addlinespace
Image & SynthID-Image & DNN-based & Multi-bit & Low & $>$99\% common edits \\
\addlinespace
Image & StegaStamp & Encoder-decoder & $\sim$100 bit & Low & Print-and-capture \\
\bottomrule
\end{tabular}
\end{table}


%% ------------------------------------------------------------
\subsection{Alignment, Defense, and Agentic Security}
\label{sec:alignment}
%% ------------------------------------------------------------

\subsubsection{Alignment Techniques}
\label{sec:alignment-tech}

Modern alignment begins with RLHF~\cite{ouyang2022rlhf}---supervised fine-tuning, reward model training, and PPO optimization---demonstrating that alignment can substitute for scale (1.3B InstructGPT preferred over 175B GPT-3). Direct Preference Optimization (DPO)~\cite{rafailov2023dpo} eliminates the separate reward model via closed-form reparameterization, becoming the dominant method by 2024. Constitutional AI~\cite{bai2022constitutionalai} replaces human labelers with AI feedback guided by a written constitution (RLAIF), reducing annotator dependence.

However, alignment remains fragile. \emph{Reward hacking} is theoretically inevitable~\cite{skalse2022rewardhacking} with quantified scaling laws~\cite{gao2023scalinglaws}. Safety training is \emph{superficial}: 10 adversarial fine-tuning examples at \$0.20 remove GPT-3.5 Turbo's guardrails~\cite{qi2024finetuning}; refusal is mediated by a single direction in activation space~\cite{arditi2024refusal}. The \emph{alignment tax}---measurable reasoning degradation from safety training---persists~\cite{huang2025safetytax}.


\subsubsection{Defense Against Adversarial Manipulation}
\label{sec:defense}

To counter the strategies categorized in Section~\ref{sec:threat-adversarial}, researchers have developed multi-layered defenses.

\paragraph{Input layer.}
Perplexity filtering detects $\sim$90\% of GCG-style adversarial suffixes but cannot catch human-crafted jailbreaks~\cite{alon2023perplexity}. SmoothLLM~\cite{robey2023smoothllm} applies randomized perturbations and majority voting, reducing ASR below 1\% across five model families.

\paragraph{System layer.}
The instruction hierarchy~\cite{wallace2024hierarchy} trains LLMs to prioritize system $>$ user $>$ tool instructions, drastically improving robustness even for unseen attacks. Instructional Segment Embeddings~\cite{wu2025ise} achieve 18.68\% improvement in robust accuracy. StruQ~\cite{chen2025struq} enforces formal prompt/data separation.

\paragraph{Output layer.}
Llama Guard~\cite{inan2023llamaguard} classifies against customizable safety taxonomies; ShieldGemma~\cite{zeng2024shieldgemma} provides +10.8\% AU-PRC improvement.

\paragraph{Model layer.}
Representation engineering~\cite{zou2023repe} extracts concept vectors for mechanistic steering. Circuit breakers~\cite{zou2024circuitbreakers} interrupt harmful generation at the representation level (ASR $\sim$3.8\% on HarmBench), though multi-turn attacks still achieve $\sim$54\%. Adversarial training via R2D2~\cite{mazeika2024harmbench} and latent adversarial training~\cite{sheshadri2024lat} (700$\times$ fewer passes) incorporate attacks during training.

\paragraph{The defense gap.}
HarmBench~\cite{mazeika2024harmbench} (510 behaviors, 18 attacks, 33 LLMs) finds no single attack/defense universally effective. StrongREJECT~\cite{souly2024strongreject} reveals jailbreaks reduce capabilities---best attack scores only 0.37/1.0 on GPT-4o. Yet Nasr et al.\ showed \textbf{all 12 recently proposed defenses bypassed} with ASR $>$90\% under adaptive attacks~\cite{nasr2025attacker}. Defense-in-depth reduces ASR from 46.5\% to 5.6\%~\cite{dong2024hierarchical}, but cannot resist simultaneous adaptive attacks.


\subsubsection{Securing Agentic AI Systems}
\label{sec:agentic-security}

%% *** BLUE OCEAN — EXPANDED (+300 words) ***

To secure against the agentic threats identified in Sections~\ref{sec:threat-adversarial} and~\ref{sec:cap-autonomy}, a new class of defenses targeting the agent--tool interaction layer is emerging.

\paragraph{(a) Architecture-level defenses.}
Classical security principles---trust boundaries, context isolation, least privilege, privilege separation---must be adapted for LLM agents. Google's multi-layer defense for Gemini operationalizes these with Agent Origin Sets (constraining browseable domains), prompt injection classifiers, User Alignment Critics (a second isolated model verifying intended actions), acknowledgment gates, and continuous red teaming~\cite{google2025geminidefense}. OpenAI's governance principles include evaluating suitability, constraining action-space, maintaining legibility (audit trails), and ensuring interruptibility~\cite{shavit2024agenticgovernance}. The emerging discipline of \emph{harness engineering}---designing constraints, feedback loops, and lifecycle management around AI agents---addresses reliability but intersects critically with security: its core components (permission boundaries, human-in-the-loop gates, tool access scoping) directly instantiate defense-in-depth principles for agentic systems.

\paragraph{(b) MCP-specific security.}
Input validation must enforce parameterized queries to prevent SQLi$\to$stored prompt injection chains. Tool call validation verifies authorized scope before each invocation. The MCP specification has evolved rapidly: no authentication (November 2024), OAuth 2.1 with PKCE (March 2025), Resource Server classification with RFC 8707 Resource Indicators (June 2025), and URL-based client registration with enhanced authorization flows (November 2025)~\cite{mcpspec2025}. CVE-2025-49596 in MCP Inspector (CVSS 9.4) demonstrated that elementary session tokens and origin verification were initially absent~\cite{oligo2025inspector}. Errico et al.\ propose five control categories: per-user authentication, provenance tracking, containerized sandboxing, inline policy enforcement, and centralized governance via private registries~\cite{errico2025mcpsecurity}.

\paragraph{(c) Monitoring and runtime verification.}
Agent behavior audit trails logging all tool calls, decisions, and reasoning are essential for forensics and risk management. Anomaly detection must identify tool-call patterns deviating from expected behavior. AgentSpec~\cite{wang2025agentspec} introduces a lightweight DSL for specifying runtime constraints, preventing $>$90\% of unsafe code executions. Pro2Guard~\cite{pro2guard2025} extends this with probabilistic model checking via DTMC learning to anticipate unsafe states before violations occur. Human-in-the-loop confirmation remains critical for high-risk operations, yet only \textbf{47\% of organizations} have implemented GenAI-specific security controls~\cite{microsoft2026cyberpulse}.

\paragraph{(d) Open challenges.}
Tool poisoning via MCP descriptions achieves $>$70\% ASR through shadowing and preference manipulation~\cite{invariant2025toolpoisoning}. Cross-agent trust chain management faces fundamental challenges: a comprehensive threat analysis identified $\sim$1,700 candidate threats across twelve risk domains for multi-agent systems, including self-replicating prompt worms, latent memory poisoning, and non-deterministic planning divergence~\cite{multiagent2026threats}. Simon Willison's ``lethal trifecta''---access to sensitive data, exposure to untrusted content, and ability to communicate externally---defines the risk boundary, but most useful agentic applications inherently require all three. Mandatory access control assumes deterministic program behavior incompatible with LLM non-determinism; no consensus exists on capability granularity for AI agents. OpenAI frames prompt injection defense as ``a long-term AI security challenge'' requiring ``frontier research''~\cite{openai2025atlas}.

\begin{figure}[htbp]
\centering
\fcolorbox{green!50}{green!5}{\parbox{0.95\linewidth}{\centering\textbf{Figure~3: Layered Defense Architecture for AIGC Security}\\[3pt]
\small Detection $\leftrightarrow$ Watermarking $\leftrightarrow$ Alignment $\leftrightarrow$ Agentic Security $\leftrightarrow$ Governance}}
\end{figure}
